\documentclass[conference]{IEEEtran}

\usepackage{amsmath}
\usepackage{booktabs}
\usepackage{microtype}
\PassOptionsToPackage{hyphens}{url}
\usepackage[hidelinks]{hyperref}
\usepackage{xspace}
\usepackage{balance}

\title{Artifact Abstract: \\
       Do Code LLMs Prefer Insecure Code?\\
       A Probabilistic Study of Security Misalignment}

\author{
  \IEEEauthorblockN{Rui Melo, Sofia Reis, Andre Catarino, Rui Abreu}
}

\begin{document}

\maketitle

\begin{abstract}
This artifact accompanies the ICST 2026 paper \emph{``Do Code LLMs Prefer
Insecure Code? A Probabilistic Study of Security Misalignment.''}
It provides the \emph{DeltaSecommits} dataset of real-world vulnerable/safe
code pairs, pre-computed probabilistic alignment scores for six code LLMs,
and all analysis scripts and figures from the paper.
The artifact is packaged as a Docker image and publicly hosted on GitHub,
Docker Hub, and Hugging Face. \textbf{No GPU or internet access is needed
to reproduce the main results}---all pre-computed model outputs are included.
\end{abstract}

% ---- 1. Overview ----
\vspace{0.5cm}
\section{Overview}

Security misalignment in code LLMs is a subtle yet consequential risk: a model
that assigns higher likelihood to vulnerable code will tend to generate and rank
insecure code higher in practice. This artifact operationalises security
alignment as a \emph{preference problem} inspired by Direct Preference
Optimisation (DPO). For each pair
$(\textit{chosen}, \textit{rejected}) = (\textit{safe}, \textit{vulnerable})$
derived from real-world security commits, we measure whether the model assigns
higher log-probability to the safe variant. Three complementary signals are
computed per pair and per model:

\begin{itemize}
  \item \textbf{Log-probability (preference)}---which variant the model is more
        likely to generate.
  \item \textbf{Perplexity (fluency)}---how natural each code snippet appears
        to the model; $\Delta\text{PPL} > 0$ means the model finds vulnerable
        code more surprising.
  \item \textbf{Shannon entropy (confidence)}---mean token-level uncertainty;
        lower entropy reflects more confident predictions.
\end{itemize}

\noindent Statistical significance is assessed via Wilcoxon signed-rank tests.

% ---- 2. Dataset ----
\vspace{0.5cm}
\section{Dataset: DeltaSecommits}

The dataset is built from \emph{Secommits}, a collection of OSV-sourced
security commits. Raw commits are fetched from the GitHub API, filtered to
single-file, single-commit patches in ten source-code languages
(\texttt{.java}, \texttt{.py}, \texttt{.c}, \texttt{.cpp}, \texttt{.js},
\texttt{.ts}, \texttt{.go}, and others), and restricted to CWEs with
$\geq 30$ samples. After deduplication by vulnerability ID, the result is a
JSONL of matched (vulnerable, safe) pairs annotated with CWE, CVSS score, and
commit metadata. Each record follows the DPO schema with
\texttt{chosen} (safe code) and \texttt{rejected} (vulnerable code):

\smallskip
{\footnotesize
\begin{verbatim}
{ "cwe": ["CWE-89"],
  "chosen": "<safe code snippet>",
  "rejected": "<vulnerable code snippet>",
  "vuln_id": "GHSA-xxxx", "score": 7.5,
  "additions": 3, "deletions": 1,
  "published_date": "2021-07-27",
  "commit_href": "https://github.com/..." }
\end{verbatim}
}

\noindent The final dataset is hosted on Hugging~Face at
\href{https://huggingface.co/datasets/rufimelo/DeltaSecommits}%
{\texttt{rufimelo/DeltaSecommits}} and can be used independently of the rest
of the artifact.
All data is sourced exclusively from public repositories and open vulnerability
databases; no PII, credentials, or proprietary code is included.

% ---- 3. Models ----
\vspace{0.5cm}
\section{Models Evaluated}

Six code-focused LLMs were evaluated in 4-bit NF4 quantisation via
\texttt{bitsandbytes} (Table~\ref{tab:models}) using forward-pass inference
only---no fine-tuning is performed. This quantisation strategy enables
inference on consumer-grade GPUs ($\geq$16\,GB VRAM) while preserving model
fidelity for the log-probability comparisons used in the study.

\begin{table}[h]
\centering
\caption{Code LLMs evaluated in the study.}
\label{tab:models}
\footnotesize
\begin{tabular}{ll}
\toprule
\textbf{Label} & \textbf{Model ID} \\
\midrule
StarCoder2-7B  & \texttt{bigcode/starcoder2-7b} \\
StarCoder2-3B  & \texttt{bigcode/starcoder2-3b} \\
Mellum-4B      & \texttt{JetBrains/Mellum-4b-base} \\
DeepSeek-Coder & \texttt{deepseek-ai/} \\
               & \quad\texttt{deepseek-coder-6.7b-base} \\
CodeLlama-7B   & \texttt{meta-llama/CodeLlama-7b-hf} \\
CodeLlama-13B  & \texttt{meta-llama/CodeLlama-13b-hf} \\
\bottomrule
\end{tabular}
\end{table}

The framework is extensible: additional models can be evaluated by updating
\texttt{security\_alignment\_config.yaml} and re-running Step~3.

% ---- 4. Metrics ----
\vspace{0.5cm}
\section{Metrics}

Table~\ref{tab:metrics} summarises the six quantitative metrics recorded for
every (safe, vulnerable) code pair and model. Together they provide
complementary views of \emph{preference} (which variant the model is more
likely to generate), \emph{fluency} (how natural the model finds each variant),
\emph{confidence} (how certain the model is at the token level), and
\emph{alignment loss} (a scalar derived from the DPO objective).

\begin{table}[h]
\centering
\caption{Metrics computed per sample and model.}
\label{tab:metrics}
\footnotesize
\begin{tabular}{p{1.8cm}p{5.5cm}}
\toprule
\textbf{Metric} & \textbf{Description} \\
\midrule
Log-probability & Sum of token log-probs; higher value means the model
                  prefers the snippet. \\
Aligned         & \texttt{True} if $\log p(\text{safe}) > \log p(\text{vuln})$. \\
Perplexity      & Exponentiated average negative log-likelihood; lower
                  means more fluent. \\
PPL diff        & $\text{ppl}(\text{vuln}) - \text{ppl}(\text{safe})$;
                  positive means the model finds vulnerable code more surprising. \\
Uncertainty     & Mean Shannon entropy over token distributions; lower means
                  more confident predictions. \\
DPO loss        & $\text{softplus}(-\beta \cdot (\log p(\text{safe}) -
                  \log p(\text{vuln})))$; lower means better alignment. \\
\bottomrule
\end{tabular}
\end{table}

% ---- 5. Artifact Contents ----
\vspace{0.5cm}
\section{Artifact Contents}

The repository ($\sim$130\,MB total, all artifacts pre-included) contains:

\begin{itemize}
  \item \texttt{secommits-raw.json} ($\sim$50\,MB)---raw OSV-sourced dataset.
  \item \texttt{secommits\_filtered\_final.jsonl} ($\sim$15\,MB)---cleaned pairs.
  \item \texttt{security\_alignment/data/CWE-*.jsonl} ($\sim$5\,MB)---per-CWE
        DPO-style alignment datasets (chosen/rejected pairs).
  \item \texttt{security\_alignment/raw\_data.csv} ($\sim$10\,MB)---flat CSV
        of all per-sample model scores.
  \item \texttt{security\_alignment/*\_results/} ($\sim$30\,MB total)---per-model
        result files for all six models (log-prob, PPL, entropy, DPO loss).
  \item \texttt{security\_alignment/all\_models\_results/} ($\sim$5\,MB)---merged
        results across all models.
  \item \texttt{artifacts/plots/} ($\sim$10\,MB)---publication-ready PDF figures
        including alignment graphs, perplexity/entropy comparisons,
        DPO heatmaps, and Wilcoxon test summaries.
  \item \texttt{analysis.py} / \texttt{analysis.ipynb}---scripts to regenerate
        all paper figures from pre-computed results.
  \item \texttt{validate.py}---smoke-test: checks imports, artifact integrity,
        and Step~2/4 logic without GPU or internet access.
\end{itemize}

% ---- 6. Pipeline ----
\vspace{0.5cm}
\section{Analysis Pipeline}

The artifact implements a four-step pipeline. \textbf{Steps 1--3 require a
GPU and internet access; all intermediate outputs are pre-included so
evaluators can skip ahead to Step~4 or directly to the analysis.}

\paragraph{Step 1 -- Dataset processing.}
\texttt{sec\_aware\_cl/secommits/process\_json.py} fetches commit diffs from
the GitHub API, applies language and CWE filters, and writes
\texttt{secommits\_filtered\_final.jsonl}.

\paragraph{Step 2 -- Alignment dataset construction.}
\texttt{dataset\_builder.py} converts the filtered JSONL into per-CWE
DPO-style files under \texttt{security\_alignment/data/}.

\paragraph{Step 3 -- Forward-pass inference.}
\texttt{security\_alignment.py} loads each quantised model and runs forward
passes on every code pair, writing per-CWE JSONL results (scores, flags,
DPO loss) and a summary \texttt{alignment\_stats.jsonl} per model. The target
model list is controlled via \texttt{security\_alignment\_config.yaml}.

\paragraph{Step 4 -- Result merging.}
\texttt{join\_results.py} combines the six per-model result directories into
\texttt{all\_models\_results/} for cross-model analysis.

% ---- 7. Reproducibility ----
\vspace{0.5cm}
\section{Reproducibility}

\paragraph{Assumed reviewer skills.}
Basic command-line proficiency and Docker are sufficient to reproduce all paper
figures. Familiarity with Python/conda and a Hugging~Face account are needed
only for the full inference pipeline (Step~3). A CUDA GPU ($\geq$16\,GB VRAM)
is required only to re-run model inference.

% \begin{table}[ht]
% \centering
% \caption{Hardware requirements by task.}
% \label{tab:hw}
% \footnotesize
% \begin{tabular}{p{3.0cm}p{2.0cm}p{1.8cm}}
% \toprule
% \textbf{Task} & \textbf{GPU} & \textbf{Disk} \\
% \midrule
% Reproduce figures (Steps~4 + analysis) & Not required & $\sim$130\,MB \\
% Re-run inference (Step~3) & $\geq$16\,GB VRAM & 130\,MB + weights \\
% \bottomrule
% \end{tabular}
% \end{table}

\paragraph{Recommended path for evaluators (Docker).}
Pull the image and reproduce all paper figures from pre-computed results
(omit \texttt{--gpus all} if no GPU is available):

{\small\begin{verbatim}
docker pull rufimelo/lm-insecure-bias:latest
docker run -it --rm --gpus all \
  -v $(pwd)/artifacts:/workspace/artifacts \
  -e GITHUB_BEARER_TOKEN=$GITHUB_BEARER_TOKEN \
  -e HF_TOKEN=$HF_TOKEN \
  rufimelo/lm-insecure-bias:latest bash
# inside the container:
python validate.py
python sec_aware_cl/alignment/analysis.py
\end{verbatim}}


\paragraph{Full pipeline (GPU required).}
Reviewers with a CUDA GPU may rerun the four pipeline steps for end-to-end
validation. Full instructions are in the repository README.

% ---- 8. Availability ----
\vspace{0.5cm}
\section{Availability and Badges}

The artifact is publicly and permanently available at:

\begin{itemize}
  \item \textbf{GitHub}: {\footnotesize\url{https://github.com/rufimelo99/lm-insecure-bias}}
  \item \textbf{Docker Hub}: {\small\texttt{rufimelo/lm-insecure-bias:latest}}
  \item \textbf{Dataset (HF)}:
        \href{https://huggingface.co/datasets/rufimelo/DeltaSecommits}%
             {\texttt{rufimelo/DeltaSecommits}}
  \item \textbf{Software Heritage (Permalink)}:
        \href{https://archive.softwareheritage.org/swh:1:dir:7d83a351c04d394cd0800bfd1662ecc813245825;origin=https://github.com/rufimelo99/lm-insecure-bias;visit=swh:1:snp:c4cd318e8d1e3bbce65b7ffa9af54215c621b92b;anchor=swh:1:rev:f1c9ca390901276f2e42a397539f30e12ce2c549}%
             {\texttt{lm-insecure-bias}}
   \item \textbf{Software Heritage (SWH ID)}:
            \url{swh:1:dir:7d83a351c04d394cd0800bfd1662ecc813245825}
\end{itemize}

\noindent\textbf{Artifact badges.}
\emph{Artifact Available}: all components are openly hosted with no access
restrictions and archived on Software Heritage for long-term preservation.
\emph{Artifact Reviewed}: the artifact is functional (\texttt{validate.py}
smoke-test), documented (README + Docker), and reproducible (all paper figures
regenerable from pre-computed results without GPU access).

\balance
\end{document}
